\section{How a Liturgical Attachment Became Proper Law}

The 1988 decrees first defined the patrimony through acts. The 10 September instrument granted use of four categories of books in force in 1962: Missal, Ritual, Pontifical, and Roman Breviary. The 18 October erection decree connected priestly sanctification to the Mass and observance of liturgical and disciplinary traditions, extended use to members and certain priests in FSSP houses or churches, and required legal observance in sacramental ministry.

The constitutions made that inherited arrangement an internal norm. The excerpt published by the FSSP identifies the society as clerical and of pontifical right; places the Sacrifice of the Mass at the center of spirituality; professes fidelity to the Roman Pontiff; names sanctification of priests through priestly ministry as the general aim; and calls faithful observance of the traditions described in \emph{Ecclesia Dei} the particular means. Priestly formation, retreats, parish work accepted from bishops, education, and auxiliary vocations follow from that center.

``Patrimony'' therefore means more than permission. In the canonical analysis of Dawid Pietras, an FSSP priest and canonist, the founding acts, constitutions, and directories join older liturgy, communion with Peter, common life, Thomistic formation, and priestly apostolate. His argument is useful because it makes the internal legal case explicit. Its conclusion remains that of a knowledgeable member and scholar, not a decree resolving every later conflict.

\subsection{Rightful autonomy and its boundaries}

Canons applied through canon 732 recognize rightful autonomy so that a society can maintain its discipline and patrimony. Autonomy is neither independence nor immunity. The Holy See approves the fundamental code and supervises a pontifical-right society; local ordinaries retain their sphere; members remain bound by universal clerical law.

That framework explains why the FSSP could make the older books normative for its ordinary corporate life while a Roman authority could still ask whether a member possessed a universal faculty to celebrate according to the reformed books. The dispute is not accurately mapped as ``tradition versus obedience.'' It concerns two real legal goods: protection of the approved specificity holding the society together, and the authority of universal and superior law within the Roman Rite.

\subsection{Not a rejection clause}

The public constitutional excerpt does not define the FSSP by denial of the reformed sacraments or Vatican II. It explicitly places the foundation in the spirit of \emph{Ecclesia Dei}, professes fidelity to the pope, and cites \emph{Lumen gentium}, \emph{Presbyterorum ordinis}, and the Holy See's program of priestly formation. This differs from saying every member holds the same theological evaluation of every postconciliar development. It does show that institutional fidelity and conciliar reference are present in the approved self-description.

John Paul II's 1998 address for the tenth anniversary of \emph{Ecclesia Dei} articulated the broader settlement: legitimate diversity should work toward unity; the postconciliar liturgical reform remained valid; attachment to earlier forms deserved understanding; and bishops were to provide pastoral attention. The FSSP lived inside that double affirmation, not on the assumption that one half had silently canceled the other.

\claimcheck{``The FSSP's charism is only permission to say the 1962 Mass.''}{The liturgical faculty is central but narrower than the institutional patrimony. The approved public texts join worship to priestly sanctification, formation, common life, fidelity to the Roman Pontiff, and apostolate under bishops. A faculty concerns acts; a patrimony organizes a society's life.}
